\documentclass[11pt]{article}

\usepackage{amsmath, amssymb, amsthm, mathtools}

\newtheorem{theorem}{Theorem}
\newtheorem{lemma}{Lemma}
\newtheorem{corollary}{Corollary}
\newtheorem{definition}{Definition}
\newtheorem{remark}{Remark}

\DeclareMathOperator{\homdens}{t}

\begin{document}

\section*{Even girth implies infinitely many tensor-product counterexamples}

Throughout, graphs are finite and simple.  For a graph \(H\), write
\[
    v(H)=|V(H)|,\qquad e(H)=|E(H)|.
\]
Let \(P_H(q)\) denote the chromatic polynomial of \(H\).

\begin{definition}[Turán graphon]
For an integer \(q\geq 2\), let \(T_q\) denote the balanced complete
\(q\)-partite graphon.  Equivalently, partition \([0,1]\) into \(q\)
equal intervals \(I_1,\dots,I_q\), and define
\[
    T_q(x,y)=
    \begin{cases}
        0, & x,y\in I_i \text{ for some } i,\\
        1, & x\in I_i,\ y\in I_j,\ i\neq j.
    \end{cases}
\]
Then
\[
    e(T_q)=1-\frac1q.
\]
\end{definition}

For a graphon \(W\), write \(t(H,W)\) for the homomorphism density of
\(H\) in \(W\).  For \(T_q\), the homomorphism density of \(H\) is
\[
    t(H,T_q)=\frac{P_H(q)}{q^{v(H)}}.
\]
Indeed, a random map \(V(H)\to [q]\) contributes exactly when it is a
proper \(q\)-colouring of \(H\).

Define the normalised ratio
\[
    F_H(q)
    :=
    \frac{t(H,T_q)}{e(T_q)^{e(H)}}
    =
    \frac{P_H(q)}{q^{v(H)}}
    \left(1-\frac1q\right)^{-e(H)}.
\]
Equivalently,
\[
    t(H,T_q)
    =
    \left(1-\frac1q\right)^{e(H)} F_H(q).
\]
The value \(F_H(q)\) compares the \(H\)-density in the Turán graphon
\(T_q\) with the \(H\)-density in an Erdős--Rényi graphon of the same
edge density \(1-1/q\).

\begin{lemma}[First correction term of \(F_H(q)\)]
Let \(H\) have finite girth \(g\), and let \(c_g(H)\) be the number of
cycles of length \(g\) in \(H\).  Then, as \(q\to\infty\),
\[
    F_H(q)
    =
    1
    +
    (-1)^g \frac{c_g(H)}{(q-1)^{g-1}}
    +
    O_H\!\left((q-1)^{-g}\right).
\]
In particular, if \(g\) is even, then
\[
    F_H(q)>1
\]
for all sufficiently large \(q\).
\end{lemma}

\begin{proof}
We use the Whitney expansion of the chromatic polynomial:
\[
    P_H(q)
    =
    \sum_{A\subseteq E(H)}
    (-1)^{|A|} q^{c(A)},
\]
where \(c(A)\) is the number of connected components of the spanning
subgraph \((V(H),A)\).  Let
\[
    r(A):=v(H)-c(A)
\]
be the rank of \(A\).  Then
\[
    P_H(q)
    =
    \sum_{A\subseteq E(H)}
    (-1)^{|A|} q^{v(H)-r(A)}.
\]

We determine the first coefficients of \(P_H(q)\).  Since \(H\) has
girth \(g\), every set of fewer than \(g\) edges is acyclic.  Hence, for
\(0\leq j\leq g-2\), the sets \(A\subseteq E(H)\) with rank \(j\) are
precisely the \(j\)-edge forests, and there are \(\binom{e(H)}{j}\) of
them.  Therefore the coefficient of \(q^{v(H)-j}\) is
\[
    (-1)^j \binom{e(H)}{j},
    \qquad 0\leq j\leq g-2.
\]

Now consider rank \(g-1\).  Every \((g-1)\)-edge subset is still a
forest, so these contribute
\[
    (-1)^{g-1}\binom{e(H)}{g-1}.
\]
The only cyclic edge sets of rank \(g-1\) are the \(g\)-cycles
themselves.  Indeed, a cyclic connected component of rank \(g-1\) has
exactly \(g\) vertices.  Since the girth is \(g\), it must contain a
cycle using all \(g\) vertices.  Any additional edge would be a chord and
would create a shorter cycle, which is impossible.  Hence the cyclic
rank-\((g-1)\) sets are exactly the \(c_g(H)\) cycles of length \(g\),
and they contribute
\[
    (-1)^g c_g(H).
\]
Thus the coefficient of \(q^{v(H)-g+1}\) is
\[
    (-1)^{g-1}\binom{e(H)}{g-1}
    +
    (-1)^g c_g(H).
\]

On the other hand,
\[
    q^{v(H)}
    \left(1-\frac1q\right)^{e(H)}
    =
    q^{v(H)}
    \sum_{j=0}^{e(H)}
    (-1)^j \binom{e(H)}{j} q^{-j}.
\]
Comparing the two expansions gives
\[
    P_H(q)
    =
    q^{v(H)}
    \left(1-\frac1q\right)^{e(H)}
    +
    (-1)^g c_g(H) q^{v(H)-g+1}
    +
    O_H(q^{v(H)-g}).
\]
Dividing by
\[
    q^{v(H)}
    \left(1-\frac1q\right)^{e(H)}
\]
yields
\[
    F_H(q)
    =
    1
    +
    (-1)^g c_g(H) q^{-(g-1)}
    \left(1-\frac1q\right)^{-e(H)}
    +
    O_H(q^{-g}).
\]
Since
\[
    \left(1-\frac1q\right)^{-e(H)}
    =
    1+O_H(q^{-1}),
\]
we obtain
\[
    F_H(q)
    =
    1
    +
    (-1)^g c_g(H) q^{-(g-1)}
    +
    O_H(q^{-g}).
\]
Finally,
\[
    q^{-(g-1)}
    =
    (q-1)^{-(g-1)}
    +
    O_g\!\left((q-1)^{-g}\right),
\]
so
\[
    F_H(q)
    =
    1
    +
    (-1)^g \frac{c_g(H)}{(q-1)^{g-1}}
    +
    O_H\!\left((q-1)^{-g}\right).
\]
If \(g\) is even, then \(c_g(H)>0\) and the leading correction term is
positive.  Hence \(F_H(q)>1\) for all sufficiently large \(q\).
\end{proof}

We shall use the tensor product of graphons.  If \(U\) and \(W\) are
graphons, their tensor product is the graphon
\[
    (U\otimes W)((x_1,x_2),(y_1,y_2))
    :=
    U(x_1,y_1)W(x_2,y_2)
\]
on the product probability space.  It satisfies
\[
    e(U\otimes W)=e(U)e(W),
\]
and, for every finite graph \(H\),
\[
    t(H,U\otimes W)=t(H,U)t(H,W).
\]
Iterating, if \(W=\bigotimes_i W_i\), then
\[
    e(W)=\prod_i e(W_i),
    \qquad
    t(H,W)=\prod_i t(H,W_i).
\]

\begin{lemma}[Telescoping tensor products]
Fix \(k\geq 2\).  For each integer \(Q\geq 1\), define
\[
    W_{k,Q}
    :=
    \bigotimes_{j=1}^{Q}
    T_{Q(k-1)+j}.
\]
Then
\[
    e(W_{k,Q})=1-\frac1k.
\]
Moreover,
\[
    \lim_{Q\to\infty} t(H,W_{k,Q})
    =
    \left(1-\frac1k\right)^{e(H)}.
\]
\end{lemma}

\begin{proof}
First, by multiplicativity of edge density under tensor products,
\[
    e(W_{k,Q})
    =
    \prod_{j=1}^{Q}
    \left(1-\frac{1}{Q(k-1)+j}\right).
\]
The product telescopes:
\[
    \prod_{j=1}^{Q}
    \left(1-\frac{1}{Q(k-1)+j}\right)
    =
    \prod_{n=Q(k-1)+1}^{Qk}
    \frac{n-1}{n}
    =
    \frac{Q(k-1)}{Qk}
    =
    1-\frac1k.
\]

Next, using
\[
    t(H,T_m)
    =
    \left(1-\frac1m\right)^{e(H)}F_H(m),
\]
we get
\[
\begin{aligned}
    t(H,W_{k,Q})
    &=
    \prod_{j=1}^{Q}
    t\!\left(H,T_{Q(k-1)+j}\right)\\
    &=
    \prod_{j=1}^{Q}
    \left[
        \left(
            1-\frac{1}{Q(k-1)+j}
        \right)^{e(H)}
        F_H(Q(k-1)+j)
    \right]\\
    &=
    \left[
        \prod_{j=1}^{Q}
        \left(
            1-\frac{1}{Q(k-1)+j}
        \right)
    \right]^{e(H)}
    \prod_{j=1}^{Q} F_H(Q(k-1)+j)\\
    &=
    \left(1-\frac1k\right)^{e(H)}
    \prod_{j=1}^{Q} F_H(Q(k-1)+j).
\end{aligned}
\]

It remains to show that
\[
    \prod_{j=1}^{Q} F_H(Q(k-1)+j)\to 1.
\]
Since \(H\) has even girth, \(g\geq 4\).  From the previous lemma,
\[
    F_H(m)=1+O_H(m^{-(g-1)}).
\]
Therefore, for fixed \(k\),
\[
    \sum_{j=1}^{Q}
    \left|
        F_H(Q(k-1)+j)-1
    \right|
    \leq
    C_H
    \sum_{j=1}^{Q}
    (Q(k-1)+j)^{-(g-1)}
    \leq
    C_{H,k} Q^{2-g}.
\]
Because \(g\geq 4\), we have \(Q^{2-g}\to 0\).  Hence
\[
    \sum_{j=1}^{Q}
    \left|
        F_H(Q(k-1)+j)-1
    \right|
    \to 0.
\]
This implies
\[
    \prod_{j=1}^{Q} F_H(Q(k-1)+j)\to 1.
\]
Consequently,
\[
    \lim_{Q\to\infty} t(H,W_{k,Q})
    =
    \left(1-\frac1k\right)^{e(H)}.
\]
\end{proof}

\begin{theorem}[Even girth gives infinitely many tensor counterexamples]
Let \(H\) be a finite simple graph of even girth \(g\).  Then there
exists \(k_0=k_0(H)\) such that, for every integer \(k\geq k_0\), the
Turán graphon \(T_k\) is not optimal among tensor products of Turán
graphons at edge density \(1-1/k\).

More precisely, for every \(k\geq k_0\), there exists \(Q\geq 1\) such
that
\[
    e(W_{k,Q})=e(T_k)=1-\frac1k
\]
and
\[
    t(H,W_{k,Q})<t(H,T_k),
\]
where
\[
    W_{k,Q}
    =
    \bigotimes_{j=1}^{Q}
    T_{Q(k-1)+j}.
\]
In particular, \(T_k\) is beaten at infinitely many densities of the
form \(1-1/k\).
\end{theorem}

\begin{proof}
By the first lemma, since \(g\) is even,
\[
    F_H(k)>1
\]
for all sufficiently large \(k\).  Choose \(k_0\) such that this holds
for all \(k\geq k_0\).

Fix \(k\geq k_0\).  Then
\[
    t(H,T_k)
    =
    \left(1-\frac1k\right)^{e(H)}F_H(k)
    >
    \left(1-\frac1k\right)^{e(H)}.
\]
On the other hand, by the telescoping tensor-product lemma,
\[
    e(W_{k,Q})=1-\frac1k
\]
for every \(Q\), and
\[
    t(H,W_{k,Q})
    \to
    \left(1-\frac1k\right)^{e(H)}
    \qquad
    \text{as } Q\to\infty.
\]
Since
\[
    \left(1-\frac1k\right)^{e(H)}
    <
    t(H,T_k),
\]
we may choose \(Q\) large enough that
\[
    t(H,W_{k,Q})<t(H,T_k).
\]
Thus \(W_{k,Q}\) is a tensor product of Turán graphons with the same edge
density as \(T_k\), but with strictly smaller \(H\)-density.

Since this argument works for every \(k\geq k_0\), there are infinitely
many such counterexamples.
\end{proof}

\begin{corollary}
If \(H\) is non-bipartite and has even girth, then \(H\) gives
infinitely many counterexamples to the optimality of Turán graphons
under the edge-density constraints
\[
    e(W)=1-\frac1k.
\]
Indeed, the previous theorem applies directly.  The non-bipartite
assumption is not needed for the proof itself; it only ensures that the
example belongs to the class of non-bipartite graphs.
\end{corollary}

\end{document}